\documentclass{article}
\usepackage[utf8]{inputenc}
\usepackage{geometry}
\usepackage{enumitem}
\usepackage{hyperref}
\usepackage{xcolor}
\usepackage{graphicx}
\usepackage{array}
\usepackage{multirow}
\usepackage{amsmath, amssymb}
\usepackage{chemfig}
\usepackage{mhchem}
\usepackage{tikz}
\usepackage{setspace}
\usepackage{soul}
\usepackage{mathtools}
\usepackage{booktabs}
\usepackage{chemformula}
\usepackage{siunitx}
\usepackage{longtable}
\geometry{margin=1in}
\setlength{\parskip}{0.5em}

\begin{document}

\usetikzlibrary{positioning, shapes.geometric, arrows.meta, decorations.pathreplacing, fit, backgrounds}
\begin{center}
    {\LARGE \textbf{SI Session Planning Form}}
\end{center}

\vspace{0.5cm}

\noindent
\begin{flushleft}
\textbf{SI Leader:} Brandon Cobb \\
\textbf{Session Week:} 5 \\
\textbf{Session \#:} 5 \\
\textbf{Instructor:} Professor Dudgeon \\
\textbf{Old Plans:} \href{https://macombedu.sharepoint.com/:b:/s/ASC-SupplementalInstructionEmbeddedTutoring/IQBvRyfNvxGGQYObfgZOKn8UARbteAbUZFaLhS7nR4OoDmQ?e=Ic6HHf}{SharePoint} \\
\textbf{Session Goal:} Students will be able to define the structure of an atom, identify how many protons, neutrons and electrons are in elements and isotopes. They will be familiar with atomic masses (the average atomic mass of element isotopes based on abundance. They will understand the intentional design of the periodic table, along with periodic trends and groups. They will be able to distinguish a metal from a nonmetal, along with the representative elements. They will know orbital diagrams, along with the rules for how they are filled (Aufbau, Hund's and the Pauli exclusion principle). This includes shells and subshells. We will discuss Noble gas configurations. They will begin to understand the molecule, how bonds are formed via valence electrons, which types of bonds are formed and lewis structures. They may touch on nomenclature and memorization of polyatomic ions. They will differentiate between molecular compounds and ionic compounds. \\ 

\vspace{1em}
\textbf{\underline{Opener}}
{\centering\large\textbf{\hl{Take Attendance}}\par}\vspace{-0.25\baselineskip}


\begin{tabular}{|p{4.2cm}|p{9.8cm}|}
\hline
\textbf{Collaborative Learning Techniques} & Group Discussion \\
\hline
\textbf{Learning Strategy} & Individual Presentation \\
\hline
\textbf{Time} & 12:00---12:10 \\
\hline
\textbf{Materials} & Ready to speak \\
\hline
\textbf{Lecture/Textbook/Old Plans/Slide References} & Computer, class roster, printed incomplete periodic tables, periodic table key \\
\hline
\end{tabular}

\vspace{0.5cm}
\noindent
{\large \textbf{Instructions (please provide a \hl{DETAILED} activity breakdown below (and/or attached}} \\
\newpage

\begin{tabular}{|p{0.9\linewidth}|}
\hline

\vspace{0.2cm}
\textit{Purpose: The periodic table is your friend in chemistry and getting used to working with it is important}

\vspace{1em}

\textbf{Overview}
You will annotate a periodic table with missing fields using \emph{every piece of chemistry knowledge you currently have}. This activity emphasizes connections between atomic structure, periodic trends, element types, electron configurations, atomic masses and atomic numbers.

\begin{enumerate}
\item Label the missing \textbf{groups} (vertical columns) and \textbf{periods} (horizontal rows). Identify the valence electron number for the representative element groups.

\item In the bottom left, match the color of the element to the \textbf{group names} where applicable (alkali metals, alkaline earth metals, halogens, noble gases, etc.).

\item Identify and outline:
\begin{itemize}
\item Representative elements
\item Metals
\item Nonmetals
\item Metalloids
\end{itemize}

\item For each of the atoms with red outlines, reference the in-class periodic table and fill in the missing fields:
\begin{itemize}
\item Atomic number
\item Atomic mass
\item Number of electrons
\item Element name
\item Atomic symbol
\end{itemize}

\item Additionally, try to identify these fields for a few of the red outlined elements:
\begin{itemize}
\item Lewis dot symbol
\item Noble gas electron configuration
\item Full electron configuration
\end{itemize}

\item For each of the subshells with red outlines, specify the shell number and shell identity (s, p, d, f):

\item Of the red elements, guess the most abundant isotope number and neutron counts

\item Add any additional relationships or facts you know (ions, polyatmic ions, presence in the world, etc.)
\end{enumerate} \\

\hline
\end{tabular}

\begin{tabular}{|p{0.9\linewidth}|}
\hline

\subsection*{Expectation}

By the end, your periodic table should visually represent:\begin{itemize}
    \item Groups (vertical columns labeled)
    \item Periods (horizontal rows labeled)
    \item Group name s(alkali metals, alkaline earth metals, halogens, noble gases, etc.)
    \item Noble gas electron configuration (example: [Ne]3s$^1$)
    \item Lewis dot symbol
    \item Atomic mass
    \item Atomic number
    \item Total electron count
    \item Electron shells
    \item Subshell notation (s, p, d, f)
    \item Atomic symbol
    \item Element name
    \item Representative element identification
    \item Metal identification
    \item Nonmetal identification
    \item Full electron configuration
\end{itemize}

Use arrows, colors, labels, and notes freely. The goal is \textbf{maximum meaningful annotation}, not neatness. \\
\hline
\end{tabular}
\newpage









% \begin{tabular}{|p{0.9\linewidth}|}
% \hline
% \begin{tikzpicture}[remember picture,overlay]
% \node at (current page.center) {\begin{tikzpicture}[scale=1]
% \node[draw, very thick, minimum width=6in, minimum height=9in] (frame) {};
% \node[scale=20, anchor=center] (N) at (-2.5,0) {\textbf{N}};
% \node[scale=20, anchor=center] (slash) at (0,0) {\textbf{/}};
% \node[scale=20, anchor=center] (A) at (2.5,0) {\textbf{A}};
% \draw[thick, dashed] (N.east) -- (slash.west);
% \draw[thick, dashed] (slash.east) -- (A.west);
% \node[rotate=90, scale=1.5] at (-5,0) {Not};
% \node[rotate=90, scale=1.5] at (5,0) {Applicable};

% \draw[dotted, thick] (-4,-3) -- (-1,-1);
% \draw[dotted, thick] (1,-1) -- (4,-3);
% \node[scale=1.2] at (-4.2,-3.3) {Unavailable};
% \node[scale=1.2] at (4.2,-3.3) {Undefined};
% \end{tikzpicture}};
% \end{tikzpicture} \\
% \hline
% \end{tabular}





\newpage
\noindent
{\large \textbf{Checking for Understanding (questions \& answers)}} \\
\begin{tabular}{|p{0.9\linewidth}|}
\hline
\begin{enumerate}
    \item How can you use what you learned about the periodic table when doing your homework?
    \item How might this skill help you take more effective notes in class?
    \item How could understanding element organization help you prepare for a test?
    \item How can you explain these concepts to a classmate during group work?
    \item What strategies would you use to check your own work using these skills?
    \item How could this knowledge help you participate in discussions or answer questions in class?
    \item How can you use patterns in the periodic table to make predictions in future assignments?
    \item What challenges did you face applying these skills, and how would you overcome them next time?
    \item How could these skills help you collaborate with peers on lab activities or projects?
    \item How might you connect what you learned today to other chemistry topics or real-world examples?
\end{enumerate} \\
\hline
\end{tabular}








\newpage
\noindent
\textbf{\underline{Activity 1}} \\
\begin{tabular}{|p{4.2cm}|p{9.8cm}|}
\hline
\textbf{Collaborative Learning Techniques} & Clusters \\
\hline
\textbf{Learning Strategy} & Mixed Medium Boardwork Model \\
\hline
\textbf{Time} & 12:10---12:35 \\
\hline
\textbf{Materials} & Expo markers, scrap paper, sharpie markers, magnets \\
\hline
\textbf{Lecture/Textbook/Old Plans/Slide References} & Chapters 1-4\\
\hline
\end{tabular}

\vspace{0.5cm}

\noindent
{\large \textbf{Instructions (please provide a \hl{DETAILED} activity breakdown below (and/or attached}} \\
\newpage

\begin{tabular}{|p{0.9\linewidth}|}
\hline
\section*{Objective}
Groups will visually represent how their understanding evolved from \textbf{Classifications of Matter} to the \textbf{Atomic Model}, showing overlaps, differences, and growth. They will use pieces of paper, drawings, and notes, and organize them on an expo board using magnets.

\section*{1. Set the Stage (5 min)}
\begin{enumerate}[label=\arabic*.]
    \item \textbf{Welcome \& Context:}
    \begin{itemize}
        \item Remind participants that the goal is to explore how their understanding has expanded over the course of the topic.
        \item Show the pre-prepared Venn diagram example on a different topic.
        \item Explain: ``You’ll use the same idea—organize your ideas to show similarities, differences, and relationships—but for the journey from Classifications of Matter $\rightarrow$ Atomic Model.''
    \end{itemize}
    \item \textbf{Explain Materials:}
    \begin{itemize}
        \item Each group receives:
        \begin{itemize}
            \item Paper pieces (plain and colored)
            \item Markers, pens, pencils
            \item Magnets
            \item Expo board or wall space
        \end{itemize}
    \end{itemize}
\end{enumerate}

\section*{2. Group Brainstorm (10--15 min)}
\begin{enumerate}[label=\arabic*.]
    \item \textbf{Prompt:}
    \begin{itemize}
        \item What you \textbf{knew} at the start (Classifications of Matter)
        \item What you \textbf{learned} along the way
        \item What you \textbf{understand now} (Atomic Model)
    \end{itemize}
    \item \textbf{Instructions:}
    \begin{itemize}
        \item Write one idea, fact, or concept per piece of paper.
        \item Draw simple visuals if helpful.
        \item Encourage creativity: use colors, symbols, or different handwriting styles.
    \end{itemize}
\end{enumerate}

\section*{3. Organize Ideas on the Board (10 min)}
\begin{enumerate}[label=\arabic*.]
    \item \textbf{Demonstrate the Format:}
    \begin{itemize}
        \item Reference the Venn diagram example:
        \begin{itemize}
            \item Center = shared ideas / overlaps
            \item Sides = unique contributions or differences
            \item Arrows or groupings = connections between ideas
        \end{itemize}
    \end{itemize}
    \item \textbf{Instructions to Groups:}
    \begin{itemize}
        \item Place pieces on the expo board using magnets.
        \item Arrange to highlight relationships, patterns, and contrasts.
        \item Groups may adopt any style they like: clusters, charts, flow arrows, thematic groupings.
    \end{itemize}
\end{enumerate} \\

\hline
\end{tabular}
\newpage


\begin{tabular}{|p{0.9\linewidth}|}
\hline
\section*{4. Present \& Reflect (10 min)}
\begin{enumerate}[label=\arabic*.]
    \item \textbf{Each Group Shares:}
    \begin{itemize}
        \item Explain their arrangement.
        \item Prompt questions:
        \begin{itemize}
            \item ``What did you notice about how your understanding changed?''
            \item ``Were there overlaps or surprises?''
            \item ``How did you decide what to cluster together?''
        \end{itemize}
    \end{itemize}
    \item \textbf{Facilitator Observations:}
    \begin{itemize}
        \item Highlight connections between groups if applicable.
        \item Reinforce growth from simple classifications $\rightarrow$ complex atomic understanding.
    \end{itemize}
\end{enumerate} \\

\hline
\end{tabular}
\newpage

\begin{tabular}{|p{0.9\linewidth}|}
\hline

\section*{5. Backup Style}
It's possible only one or two students attend and this activity would not succeed with only one student brainstorming.
If this is the case, I will role-play as a peer, providing input which is both helpful and correct in some ways, but also intentionally thought provoking with wrong ways.
Be intentional about mentioning how I will not be acting as the SI Leader and instead will behave like one of their classmates, in either a good or bad way.

\begin{enumerate}[label=\arabic*.]
    \item {Positive Reinforcement}
	\begin{itemize}
            \item Encourage students when they bring up specific content points they want to touch on.
	    \item Suggest things are a good idea when I think they are good.
	    \item Utilize the tools I brought early on to facilitate good emphasis on important points.
        \end{itemize}
    \item {Sources of Information}
        \begin{itemize}
	    \item Hand-written notes
	    \item The periodic table
	    \item Print-out notes
	    \item Lecture slides
	    \item Review materials on Canvas
	\end{itemize}
    \item {Distillation of Knowledge}
        \begin{itemize}
	    \item What question are we answering?
	    \item Does everyone agree?
	    \item What useful ideas can we adapt from our disagreements?
	    \item What tools does each person feel comfortable using?
	    \item Does anyone in the group struggle and have a weakness?
	    \item Does anyone have any strengths they can employ to help address the weaknesses?
	    \item Does anyone remember really specific concepts/diagrams they like?
	    \item What overall presentation format do we have available?
	    \item What is the ideal way everyone agrees for outr presentation?
	    \item Who wants to be the scribe? Should we all be scribing?
	    \item Who wants to present? Should we all present?
	    \item Who has the best handwriting?
	\end{itemize}
    \item {Break the dependency cycle}
        \begin{itemize}
	    \item Break the dependency cycle with wait time 1 and 2 and redirect questions.
	    \item When students push back, switch to positive reinforcement.
	    \item Try not to pull out tools like my notes or markers until later or if asked.
	    \item Ask questions instead of providing answers.
	    \item Be very clear of the instructions early on to prevent followup ambiguity.
	\end{itemize}
\end{enumerate} \\

\hline
\end{tabular}	    

% \hline
% \begin{tikzpicture}[remember picture,overlay]
% \node at (current page.center) {\begin{tikzpicture}[scale=1]
% \node[draw, very thick, minimum width=6in, minimum height=9in] (frame) {};
% \node[scale=20, anchor=center] (N) at (-2.5,0) {\textbf{N}};
% \node[scale=20, anchor=center] (slash) at (0,0) {\textbf{/}};
% \node[scale=20, anchor=center] (A) at (2.5,0) {\textbf{A}};
% \draw[thick, dashed] (N.east) -- (slash.west);
% \draw[thick, dashed] (slash.east) -- (A.west);
% \node[rotate=90, scale=1.5] at (-5,0) {Not};
% \node[rotate=90, scale=1.5] at (5,0) {Applicable};

% \draw[dotted, thick] (-4,-3) -- (-1,-1);
% \draw[dotted, thick] (1,-1) -- (4,-3);
% \node[scale=1.2] at (-4.2,-3.3) {Unavailable};
% \node[scale=1.2] at (4.2,-3.3) {Undefined};
% \end{tikzpicture}};
% \end{tikzpicture} \\
% \hline
% \end{tabular}

\newpage

\noindent
{\large \textbf{Checking for Understanding (questions)}} \\[0.2em]

\begin{tabular}{|p{0.9\linewidth}|}
\hline

\begin{enumerate}
    \item How did your group decide which ideas were most important to include on the board? How could this strategy help you in organizing future study topics?
    \item Which method of presenting your ideas (clusters, flow arrows, diagrams, colors) worked best for communicating understanding to others? Why?
    \item Were there connections or overlaps you noticed between your ideas and those of other groups? How might recognizing patterns like these help you in understanding complex concepts?
    \item How did creating visual representations (drawings, charts, Venn-style arrangements) help you see your learning progression? How could this approach be applied when reviewing other units?
    \item If you were explaining your group's board to someone who had never studied this topic, what approach would make your presentation most clear and engaging? How does this relate to presenting scientific ideas in class or on exams?
    \item Did working collaboratively help reveal gaps in your individual understanding? How can teamwork strategies be applied in future learning or lab activities?
    \item Reflect on your process from brainstorming to final presentation. What did you learn about organizing knowledge that could help you succeed in the class?
\end{enumerate} \\

\hline
\end{tabular}
\newpage

\noindent
\textbf{\underline{Activity 2}} \\
\begin{tabular}{|p{4.2cm}|p{9.8cm}|}
\hline
\textbf{Collaborative Learning Techniques} & Walkthrough \\
\hline
\textbf{Learning Strategy} & Measure, Listen, Remeasure \\
\hline
\textbf{Time} & 12:35---12:50 \\
\hline
\textbf{Materials} & A computer, textbook, ready to speak \\
\hline
\textbf{Lecture/Textbook/Old Plans/Slide References} &  Chapter 1 and 2 \\
\hline
\end{tabular}

\vspace{0.5cm}
\noindent
{\large \textbf{Instructions (please provide a \hl{DETAILED} activity breakdown below (and/or attached}} \\
\newpage
\begin{tabular}{|p{0.9\linewidth}|}
\hline
\hline
\end{tabular}
\newpage

\noindent
{\large \textbf{Content (fully worked out problems \& answers)}} \\

\begin{tabular}{|p{0.9\linewidth}|}
\hline
\hline
\end{tabular}

\newpage
\noindent
{\large \textbf{Checking for Understanding (questions \& answers)}} \\
\begin{tabular}{|p{0.9\linewidth}|}
\hline
\hline
\end{tabular}




\newpage
\noindent
\textbf{\underline{Closer}}\\
\begin{tabular}{|p{4.2cm}|p{9.8cm}|}
\hline
\textbf{Collaborative Learning Techniques} & Peer Review & Self-Reflection \\
\hline
\textbf{Learning Strategy} & Test-Taking Simulation and Error Analysis \\
\hline
\textbf{Time} & 12:50---12:55 \\
\hline
\textbf{Materials} & Previous quizzes/tests, pen/pencil, notebook, timer \\
\hline
\textbf{Lecture/Textbook/Old Plans/Slide References} & Relevant recent assessments or practice problems \\
\hline
\end{tabular}

\vspace{0.5cm}
\noindent
{\large \textbf{Instructions}} \\[0.2em]

\begin{tabular}{|p{0.9\linewidth}|}
\hline
\hline
\end{tabular}

\newpage
\noindent
{\large \textbf{Checking for Understanding (questions & answers)}} \\[0.2em]

\begin{tabular}{|p{0.9\linewidth}|}
\hline
\hline
\end{tabular}

\newpage
\textbf{Plan Checklist}

\begin{tabular}{|p{0.9\linewidth}|}
\hline
\begin{itemize}
\item[$\Box$] Variety of Collaborative Learning Techniques and Learning Strategies
\item[$\Box$] Chapter/slide \# where information is coming from
\item[$\Box$] Included timestamps (and buffer time as needed)
\item[$\Box$] Details on how leader will break up students
\item[$\Box$] Checking for Understanding questions
\item[$\Box$] Fully worked out problems and examples for each activity
\item[$\Box$] Necessary links required for online implementation (if applicable)
\item[$\Box$] Source material correctly cited (if applicable)
\item[$\Box$] Plan is uploaded to Teams
\end{itemize}
\\ \hline
\end{tabular}

\end{document}

